\section{Discussion}
\label{sec:discussion}

\subsection{Interpretation of Experimental Findings}

The experiments show that GPU-assisted AES-CTR acceleration can persist beyond isolated kernel
execution when synchronous host--device transfers and synchronization are included. The positive
prepared-buffer results on T600 and T4 identify favorable bulk-processing conditions for the
tested implementations. The T600 result also shows that the observed benefit is not confined to
the datacenter-oriented configuration. The GTX pipeline and mode-dependent outcomes establish
that resident computational performance alone does not determine whether offloading improves
host-observed processing time.

CTR permits independent counter-block encryption, provided that counter and nonce requirements
are satisfied~\cite{dworkin2001modes}. This independence supplies parallel work for the CUDA
implementation. At 32~MiB, the resident measurements show an advantage over the evaluated AES-NI
path on every system in both modes. The pipeline measurements then establish whether an
advantage remains within the broader execution boundary. Their lower throughput reflects the
additional operations in that boundary, but the measurements do not isolate how much of the
difference is attributable to physical transfers, synchronization, or runtime overhead.

The host implementation is equally relevant to interpreting acceleration. On T4, similar GPU
throughput in ECB and CTR accompanies different host-relative speedups because the corresponding
AES-NI rates differ substantially. The CTR speedup therefore characterizes the complete CPU--GPU
comparison; it does not indicate that the GPU executes CTR several times faster than ECB. The
underlying cause of the host disparity was not isolated, so attributing it to a processor
feature, library behavior, or cloud mechanism would exceed the evidence. The T600 and T4
acceleration factors also use different host processors as their denominators. Their relative
magnitudes cannot rank GPU performance independently of the host; absolute throughput and
host-relative acceleration must be considered together when comparing the tested configurations.

Workload size and session consistency further qualify these findings. Non-monotonic throughput
prevents identification of a universal CPU-to-GPU crossover threshold. A positive aggregate
speedup can also coexist with sessions favoring the CPU, as observed for T600 ECB. The results
support selecting execution paths using workload-specific measurements and paired session
outcomes, rather than peak kernel throughput or a single favorable observation. They do not
establish that 32~MiB is optimal or that larger buffers would necessarily improve performance.

\subsection{Practical Implications and Generalizability}

For telecommunications, the findings provide a basis for evaluating GPU-assisted bulk encryption
in communication gateways and server-side data-protection services. Systems already handling
substantial contiguous buffers could assess the GPU as an additional execution path alongside
AES-NI. This is a potential use of the measured capability, not a demonstrated improvement in a
deployed communication service. The relevant comparison must include the operations needed to
acquire, organize, protect, and deliver data in the intended application. Where general-purpose
GPUs are already present, this could allow bulk encryption to use existing heterogeneous
resources without a separate cryptographic accelerator. The measurements support assessing that
option against the available hardware-assisted CPU path rather than assuming that an additional
accelerator will improve the service.

Aggregating independent requests could increase the parallel work available per accelerator
invocation, but introduces requirements absent from the buffer experiments. Queueing while a
batch forms can increase latency, especially under low traffic. Packet metadata, independent
cryptographic contexts, counter allocation, and output ordering must also be maintained.
Consequently, a contiguous 32~MiB buffer is not interchangeable with a collection of
independently managed packets, and the measurements establish neither an optimal batching policy
nor a minimum packet batch size. The relevant trade-off is therefore between the processing
benefit of a larger GPU workload and the time spent waiting to assemble and manage it. Requests
with strict latency requirements may need a different execution policy from sustained bulk
traffic.

AES-CTR is specified for IPsec Encapsulating Security Payload (ESP) confidentiality in
RFC~3686~\cite{rfc3686}, while AES-GCM is specified for ESP in RFC~4106~\cite{rfc4106}. These
standards connect the evaluated encryption operation to IPsec-oriented architectures, but the
experiments cover neither complete ESP processing nor authenticated encryption. Such processing
adds packet framing, security-association state, nonce and sequence-number management,
authentication, and delivery. AES-GCM additionally requires GHASH and authentication-tag
processing. The measured CTR acceleration cannot determine the throughput of that complete
operation.

An IPsec extension must preserve counter uniqueness, authentication boundaries, and the
association between packet data and protocol metadata when organizing parallel work. Its
performance assessment would need to include both throughput and latency across the
packet-processing path. Faster counter-mode encryption could benefit that path only to the
extent that the remaining operations and their interactions permit an overall improvement.

Device-resident processing offers another possible arrangement when input data are already
available on the GPU. The resident measurements characterize encryption under this placement,
while application acquisition, authentication, and delivery remain to be evaluated. Similar
questions could arise in other data-intensive industries, such as storage or analytics services
processing large buffers. These are plausible extensions of the workload structure, not
validated application results. A workload-aware architecture retaining CPU processing for
smaller or latency-sensitive requests and considering GPU execution for bulk work remains an
untested design implication.

\subsection{Study Limitations}

Each configuration represents one host or cloud allocation with five study sessions. The
observations establish repeatability within those environments, not performance across a
population of comparable machines. Differences among the hardware and software configurations
were not independently controlled, preventing causal attribution to a single component. The
selected AES-NI library path and validated CUDA kernel also do not establish the fastest
attainable CPU or GPU implementations, including multithreaded CPU services.

The prepared-buffer boundary excludes application-level work and preparation costs specified in
Section~\ref{sec:methodology}. Its throughput is therefore not complete
telecommunications-system throughput. The fixed ascending size order confounds buffer-size
effects with possible time-dependent execution behavior; the local 16~MiB reduction cannot be
assigned to memory effects, scheduling, clocks, or thermal conditions from these measurements.
Five-session ranges and directional consistency are descriptive evidence rather than confidence
intervals or formal significance tests. These limits constrain both causal explanations and
predictions of performance beyond the tested execution paths.

\subsection{Future Research}

Three extensions follow directly from the measured acceleration and its constraints.

First, implement and evaluate an IPsec-oriented authenticated-encryption pipeline, including
AES-GCM, GHASH, tag processing, and protocol state. A representative gateway workload should
include packet preparation and delivery and compare complete throughput and latency with a
hardware-assisted CPU implementation. This would test whether the observed encryption-component
advantage survives the additional work required by the protocol.

Second, evaluate packet aggregation and workload-aware CPU--GPU selection under varying traffic
volume and latency requirements. The study should retain independent cryptographic contexts and
output ordering while measuring queueing delay as well as processing time. More densely sampled
workload sizes and additional independent systems would help assess whether a selection policy
remains effective beyond the current size grid and host allocations.

Third, compare the synchronous pipeline with overlapping transfers and computation, persistent
device-resident processing, and supported direct-data-transfer mechanisms. Measurements should
separate the relevant processing stages while also reporting the complete application interval.
This would determine whether changing data placement and transfer scheduling alters the
conditions favoring GPU execution; the present synchronous results alone do not predict those
gains.
